\documentclass[11pt,a4paper]{article}

\usepackage{amsmath,amssymb,amsthm}
\usepackage{geometry}
\geometry{margin=1in}

\title{\textbf{Solution of the Discrete Diophantine Equation for the Minkowski Interval}}
\author{\textbf{Noriaki Kihara} \\
\small\textit{Chief Architect / CEO, WF System Co., Ltd.} \\
\small\textit{Osaka University, School of Engineering Science (Alumni)} \\
\small\textit{(Pioneer in Computational Geometry and 3D CAD Architecture)}
}
\date{}

\begin{document}
\maketitle

\section*{Acknowledgments}
The author of this paper expresses gratitude to the AI technology ``AntiGravity'', which contributed to the translation of computational logic construction and conceptual paradigm shifts.

\begin{abstract}
This paper investigates the number-theoretic constraints that arise when formulating the flat Minkowski spacetime of special relativity on a discrete integer lattice $ \mathbb{Z}^{1+N} $. Assuming spatial isotropy, the invariant of the Minkowski interval reduces to the Diophantine equation
\[ l^2 + t^2 = N x^2 \]
where $ l, t, x \in \mathbb{Z}^+ $. The primary objective of this study is to determine the necessary and sufficient conditions for this equation to have non-trivial positive integer solutions. By utilizing Fermat's theorem on sums of two squares and the method of infinite descent, we prove that the necessary and sufficient condition is that in the prime factorization of $ N $, every prime factor $ p \equiv 3 \pmod{4} $ must have an even exponent. Consequently, it is shown that no non-trivial solutions exist in our 4-dimensional spacetime ($ N=3 $) (an arithmetic anomaly).

On the other hand, solutions do exist for specific dimensions known in superstring theory and M-theory, such as $ N=10 $ ($ D=11 $) and $ N=25 $ ($ D=26 $). While this paper highlights the coincidence of these dimensions from the perspective of the vanishing of conformal anomalies in continuous theories and the existence of integer solutions in discrete models, the model strictly addresses discrete number-theoretic constraints in the flat limit of special relativity. It does not intend to construct a theory of quantum gravity, propose dynamics or a Lagrangian, or provide a direct extension to continuous spacetime.

This research presents a purely arithmetic-geometric toy model, demonstrating the potential that simple Diophantine constraints on a discrete Minkowski lattice can impose strong mathematical restrictions on the selection of spacetime dimensions.
\end{abstract}

\section{Introduction}
The equation to be derived is:
\[ l^2 = x^2 + y^2 + z^2 - t^2 \]
Since inversion symmetry is assumed for \( l, x, y, z, t \), they are all assumed to be 0 or positive integer solutions.

It must be explicitly stated that this paper does not intend to construct a full-fledged physical theory such as quantum gravity or string theory. It is strictly a purely number-theoretic toy model analyzing the Diophantine equation
\[ l^2 + t^2 = N x^2 \]
that emerges when the limit of flat Minkowski spacetime in special relativity is described on a discrete integer lattice $ \mathbb{Z}^{1+N} $. This equation contains no curvature or interactions, and possesses no mechanical dynamics or Lagrangian. Under this simple setting, the main focus of this research is to investigate the conditions for the existence of positive integer solutions corresponding to isotropic geodesics.

\section{Card Game Metaphor}
Consider an infinitely vast table for a card game. First, let the dealer be \( l \).
Next, we metaphorically treat dimensions as players in the game. \\
\textbf{0-dimensional:} No player exists. \\
\textbf{1-dimensional:} The player is \( x \). \\
\textbf{2-dimensional:} The players are \( x, y \). \\
\textbf{3-dimensional:} The players are \( x, y, z \).

At the initial state (the start of the game), both the dealer and the players have 0 cards in their hands.

\paragraph{0-dimensional board:}
Since there are no players, the dealer's table \( l^2 \) has only two states: Cards do not exist ($0$) or Cards exist ($1$).

\paragraph{1-dimensional board:}
The only player is \( x \). \( x \) can receive any number of cards from the dealer, but the condition is that their hand must always be a square number. Additionally, the dealer must add exactly the requested number of cards to their own hand, which must also remain a square number. If this rule is violated, it results in a game over.
In this case, it is very simple:
\[ l^2 = x^2 \]
Thus, the number of cards in \( x \)'s hand and the dealer's hand will be the square numbers 0, 1, 4, 9. Since \( x \) can easily request a number of cards that keeps both hands as square numbers, the game never ends.

\paragraph{2-dimensional board:}
The players are \( x, y \). Assuming spatial axis symmetry, there is a special rule where \( y \) can only request the same number of cards as \( x \). That is,
\[ l^2 = x^2 + y^2 = 2x^2 \]
If \( 2x^2 \) is not a square number, it is game over. Thus, \( x^2 \) must be a multiple of 2. However, since \( x \) can ensure they always request a number that satisfies this rule, the game never ends.

\paragraph{3-dimensional board:}
The players are \( x, y, z \). Following the same logic:
\[ l^2 = x^2 + y^2 + z^2 = 3x^2 \]
If \( 3x^2 \) is not a square number, it is a game over. Thus, \( x^2 \) must be a multiple of 3. Since \( x \) can always request such a number from the dealer, the game never ends.

In other words, under these rules, on an \( n \)-dimensional board with \( n \) players:
\[ l^2 = n x^2 \]
So if \( n \) is a square number, the game never ends.

Until now, this is a rather mundane and boring game. Let us introduce a special player, the Joker \( t \). He has a unique rule and can request an arbitrary number of cards from the dealer independently of \( x \). \( t \)'s hand must also be a square number. However, the dealer must return the requested number of cards \( t \) back to the card deck from their own hand. This introduces the concept of \( -t \) as a negative. Even in this case, the deck's count must remain a square number. If it cannot be made a square number, it is a game over.

This rule can only avoid a game over when the sum of players plus the Joker (\( n \)) reaches specific numbers.
\[ l^2 = (n - 1) x^2 - t^2 \]
This condition is satisfied when:
\[ n = 1, 2, 3, 5, 6, 9, 10, 11, 14, 17, 18, 19, 21, 26, 27, 30, 33, 35, 37, \dots \]

\section{Discrete Minkowski Metric and Symmetry Assumption}
Consider the line element \( ds^2 \) of the flat Minkowski metric in \( D \)-dimensional spacetime. Let the speed of light \( c=1 \), the time axis be 1-dimensional, and the spatial axes be \( N \)-dimensional (\( D = N+1 \)).
\[ ds^2 = \sum_{i=1}^{N} (dx_i)^2 - dt^2 \]
Assume that spacetime is a discrete lattice composed of integer multiples of a minimal unit akin to the Planck length, and restrict all displacements \( dx_i, dt \) and the invariant \( ds \) to be positive integers. Here, \( ds \) will be denoted as \( l \).

To preserve spatial isotropy, we consider a symmetric geodesic where the displacement along each spatial dimension is equal in a specific inertial frame. That is, for any \( i \), let \( dx_i = x \in \mathbb{Z}^+ \). The equation is then simplified to the following form:
\[ l^2 = N x^2 - t^2 \implies l^2 + t^2 = N x^2 \]

\section{Diophantine Constraints on Dimension \( N \)}
We identify the condition on the spatial dimension \( N \) for the equation derived in the previous section to have non-trivial solutions \( (l, t, x) \in (\mathbb{Z}^+)^3 \).

\newtheorem*{theorem}{Theorem 4.1}
\begin{theorem}
A necessary and sufficient condition for the equation \( l^2 + t^2 = N x^2 \) to have positive integer solutions is that in the prime factorization of \( N \), the exponent of any prime factor \( p \equiv 3 \pmod 4 \) must be an even number.
\end{theorem}

\begin{proof}
We use proof by contradiction and Fermat's Method of Infinite Descent.

First, assume that \( N \) contains a prime factor \( p \equiv 3 \pmod 4 \) to an odd power, and set \( N = p^{2k+1}M \) (where \( p \) and \( M \) are coprime, \( k \ge 0 \)). In this case, the equation can be expressed as:
\[ l^2 + t^2 = p^{2k+1} M x^2 \]

By Fermat's theorem on sums of two squares, if a prime \( p \equiv 3 \pmod 4 \) divides the sum of two squares \( l^2 + t^2 \), then \( p \) must divide both \( l \) and \( t \).
Therefore, we can substitute \( l = p l' \) and \( t = p t' \) (\( l', t' \in \mathbb{Z}^+ \)). Substituting this into the original equation yields:
\[ p^2(l'^2 + t'^2) = p^{2k+1} M x^2 \]
Dividing both sides by \( p^2 \) gives:
\[ l'^2 + t'^2 = p^{2k-1} M x^2 \]

By repeating this operation (\( k \) times), the exponent of \( p \) on the right side ultimately becomes 1, giving the following equation:
\[ L^2 + T^2 = p M x^2 \]
Since the right side is a multiple of \( p \) again, the left side \( L, T \) must also be multiples of \( p \). Substituting \( L = p L' \) and \( T = p T' \) yields:
\[ p^2(L'^2 + T'^2) = p M x^2 \implies p(L'^2 + T'^2) = M x^2 \]
Here, since the left side is a multiple of \( p \), the right side \( M x^2 \) must also be a multiple of \( p \). However, since \( M \) is coprime to \( p \) (\( p \nmid M \)), the prime \( p \) must divide \( x^2 \), resulting in \( x \) also having \( p \) as a factor (\( x = p x' \)).

Substituting this back into the above equation gives:
\[ p(L'^2 + T'^2) = M p^2 x'^2 \implies L'^2 + T'^2 = M p x'^2 \]
Since \( N = p M \), this leads to:
\[ L'^2 + T'^2 = N x'^2 \]

This equation is of the exact same form, but its overall scale is strictly smaller (\( L' < l \), \( T' < t \), \( x' < x \)).
If a positive integer solution \( (l, t, x) \) exists, this operation would imply the capability to continually generate a set of infinitely smaller positive integer solutions. This contradicts the well-ordering principle of natural numbers (Method of Infinite Descent).
Therefore, no positive integer solution exists, restricting the solutions solely to the trivial solution \( (0, 0, 0) \).
\end{proof}

\section{Physical Interpretation}
Applying Theorem 4.1, we see that a strong constraint is placed on the spatial dimension \( N \) required for the existence of isotropic discrete geodesics.

\begin{itemize}
    \item \textbf{\( N=3 \) (\( D=4 \)):} Since \( 3 \equiv 3 \pmod{4} \), it does not satisfy the theorem's conditions, and no non-trivial positive integer solutions exist.
    \item \textbf{\( N=10 \) (\( D=11 \)):} Solutions exist (e.g., \( l=1, t=3, x=1 \)).
    \item \textbf{\( N=25 \) (\( D=26 \)):} Solutions exist (e.g., \( l=3, t=4, x=1 \)).
\end{itemize}

These results suggest that arithmetic constraints in discrete lattice models could potentially influence the selection of spacetime dimensions. However, this model remains a purely number-theoretic investigation restricted to the flat limit of special relativity.

\section{Scope and Limitations}
This research is not aimed at constructing continuous spacetime or quantum gravity theories; it is strictly limited to \textbf{simple Diophantine constraints on a discrete Minkowski lattice}.

Terms like ``correspondence'' and ``duality'' used in this paper refer only to \textbf{partial and formal similarities} between mathematical structures, and do not assert physical identity or theoretical equivalence.

What this research does NOT assert:
\begin{itemize}
    \item The construction of physical theories in continuous spacetime.
    \item The presentation of quantization methods, Lagrangians, or dynamics.
    \item Direct identity with or expansion of string theory or existing dualities.
\end{itemize}

This model is positioned merely as an example of number-theoretic constraints within the simplest discrete version of flat spacetime, void of interactions or curvature.

\section{Discussion}
The single constraint \( l^2 + t^2 = N x^2 \) (case where \( K=0 \)) addressed in this paper describes the set of integer points corresponding to the light cone (massless shell). Transformations on this constraint reduce to existing symmetries, such as the discrete Lorentz group.

If expanded more generally to \( l^2 + t^2 - N x^2 = K \) (\( K \in \mathbb{Z} \)), one could investigate the correspondence between integer points across different \( K \) values, though this paper does not perform this detailed classification.

\section{Conclusion}
In this paper, we investigated the Diophantine equation \( l^2 + t^2 = N x^2 \), which emerges when the flat Minkowski spacetime of special relativity is discretized on an integer lattice. The existence condition for its non-trivial positive integer solutions was completely determined using Fermat's theorem on sums of two squares and the method of infinite descent.

In particular, an arithmetic constraint was revealed dictating the non-existence of solutions in 4-dimensional spacetime (\( N=3 \)), while the existence of solutions for specific \( N \) (e.g., 10, 25) was also demonstrated.

Such number-theoretic constraints are intriguing as an aspect of discrete spacetime models. However, this research is strictly confined to this simple toy model and does not intend direct application to continuous theories or quantum gravity. It is our hope that this arithmetic insight will serve as a foundational reference for developing this approach in the future.

\end{document}
